\section{Scope of Price Evidence}
\label{app:scope_adaptation_submission}

The main text treats price evidence as corroboration rather
than causal identification. This appendix records the price
decomposition that keeps that interpretation disciplined: the
overlap sign reversal and the high-value-segment positive
result. The strategic-adaptation simulation is reported
separately in Appendix F.

\subsection{Overlap and Segment Scope}
\label{app:overlap_scope_submission}

The broad price association is positive, with the headline range
$\valHeadlineRange$. Under overlap-cell ATT reweighting, the
aggregate coefficient reverses to $\valScopeOverlapCoef$; under
propensity-score trimming it becomes $\valScopePSCoef$. This
does not refute the screening interpretation. It shows that the
broad estimate and the overlap estimate average different parts
of the procurement distribution.

The segment decomposition is the useful object. The lowest three
tender-value quartiles have negative coefficients
($\valSRQoneBroad$, $\valSRQtwoBroad$, and
$\valSRQthreeBroad$ in the broad specification). The highest
tender-value quartile is positive in both the broad and overlap
designs: $\valSRQfourBroad$ and $\valSRQfourATT$
($p=\valSRQfourATTp$). Trimming heavily weighted overlap cells
does not eliminate the negative aggregate estimate; it moves from
$\valSRTrimNone$ with no trim to $\valSRTrimTen$ after dropping
the top decile of cell weights and $\valSRTrimFifty$ after
dropping the top half.

The legal-economic consequence is scope, not damages. The price
results show that flagged environments are economically
non-neutral, especially in the high-value segment where cartel
rents are most plausible. They do not supply a causal price
effect, a markup, or a damages calculation.

\subsection{Selection and Within-Cell Decomposition}
\label{app:price_decomposition_submission}

The broad positive imprint and the negative overlap-cell ATT
coefficient decompose into a positive across-cell selection
component and a negative within-cell pattern. Among non-treated
items, mean log price rises monotonically across cell FL-share
quintiles (gap $\Delta=\valSelTestDelta$ log-points across cells);
a regression of non-treated log price on cell FL-share with all
five cell dimensions as marginal fixed effects yields
$\valSelTestCoefFullFE$ (SE $\valSelTestSEFullFE$,
$N=\valSelTestN$). Within overlap cells, frequent-loser presence
is associated with a winner-to-reference ratio lower by
$\valMechWinnerVsRef$ log-points (SE $\valMechWinnerVsRefSE$),
an association that becomes insignificant once
$\log(\text{n\_firms})$ is included ($\valMechWinnerVsRefControlled$,
SE $\valMechWinnerVsRefControlledSE$, n.s.); frequent-loser
presence is associated with $\valMechNFirms$ log-bidders
($\approx\valMechNFirmsPct$ more) within the same cell type. The
two associations exchange dominance at the tender-density margin:
positive in sparse tenders ($\Delta=\valMechSparseBidderDelta$ at
1--3 bidders) and negative in dense tenders
($\Delta=\valMechDenseBidderDelta$ at 11+ bidders), with the
7--10 bucket as the transition ($\valMechMidHighBidderDelta$).
These are descriptive associations consistent with economic
non-neutrality; they do not identify a mechanism, a causal price
effect, overcharges, or damages.

\begin{table}[htbp]
\centering
\caption{Selection and Within-Cell Price Decomposition of the FL-Price Sign Reversal}
\label{tab:sign_reversal_decomposition_submission}
\footnotesize
\begin{adjustbox}{max width=\textwidth,center}
\begin{threeparttable}
\begin{tabular}{p{4.8cm}cp{2.0cm}p{5.6cm}}
\toprule
Component & Estimate & SE / p & Interpretation \\
\midrule
Broad-sample baseline & $+\valBetaBroad$ & $p=\valBetaBroadP$ &
  Joint association: selection $+$ within-cell pattern \\
\midrule
\multicolumn{4}{l}{\textit{Across-cell selection (non-treated price on cell FL-share)}} \\
\hspace{6pt} Raw OLS & $+22.65$ & SE $0.03$ &
  Cells where cartels concentrate are structurally high-priced \\
\hspace{6pt} + item-group + year FE & $+24.04$ & SE $0.56$ &
  Selection beyond marginal cell dimensions \\
\hspace{6pt} + all five marginal cell-dim FE & $\valSelTestCoefFullFE$ &
  SE $\valSelTestSEFullFE$ &
  Selection coefficient after partialling out each cell dim \\
\midrule
\multicolumn{4}{l}{\textit{Within-cell price pattern (winner-to-reference ratio)}} \\
\hspace{6pt} Within cell, no controls & $\valMechWinnerVsRef$ &
  SE $\valMechWinnerVsRefSE$, $p<10^{-30}$ &
  FL presence is associated with lower winner-to-reference ratio \\
\hspace{6pt} + $\log(\text{n\_firms})$ control & $\valMechWinnerVsRefControlled$ &
  SE $\valMechWinnerVsRefControlledSE$, n.s. &
  Association is accounted for by the bidder-count channel \\
\hspace{6pt} Log-bidder inflation & $\valMechNFirms$ &
  SE $\valMechNFirmsSE$, $p<10^{-30}$ &
  $\approx \valMechNFirmsPct$ more bidders in FL-present tenders \\
\midrule
\multicolumn{4}{l}{\textit{Sign-reversal boundary by bidder count}} \\
\hspace{6pt} 1--3 bidders (sparse) & $\valMechSparseBidderDelta$ & --- &
  Selection regime: FL $=$ higher price \\
\hspace{6pt} 7--10 bidders & $\valMechMidHighBidderDelta$ & --- &
  Transition \\
\hspace{6pt} 11+ bidders (dense) & $\valMechDenseBidderDelta$ & --- &
  Within-cell regime: FL $=$ lower price \\
\bottomrule
\end{tabular}
\begin{tablenotes}
\footnotesize
\item \textit{Notes:} Selection component: log negotiated price
among non-treated items only, regressed on cell FL-share
(across-cell variation). Within-cell component: log winner-to-
reference ratio regressed on FL presence with overlap-cell fixed
effects (within-cell variation). All standard errors clustered
by overlap cell. The bidder-count rows report
$\Delta(\text{FL-present} - \text{FL-absent})$ on winner-to-
reference within each bidder-count bucket. Estimates are
descriptive associations, not causal effects, overcharges,
markups, or damages.
\end{tablenotes}
\end{threeparttable}
\end{adjustbox}
\end{table}

\subsection{Implied Monetary Scope}
\label{app:monetary_scope_submission}

We report an implied monetary-scope calculation only to indicate
institutional scale, not as a damages calculation. Combining the
price-side coefficients with the share of items in the loser-side
stratum and BEC sample-period procurement volume of
\valWelfareDenom, the descriptive scope range is
\valWelfareCrossfit{} per year (cross-fitted lower bound) to
\valWelfareIV{} (instrumented upper bound), equivalent to about
\valWelfareLowPct{}--\valWelfareHighPct{}\% of state procurement
spending in scope; the OLS midpoint is \valWelfareOLS. These
figures bound the institutional stakes the screen is reading from.
They are not damages, overcharges, or welfare losses, and they do
not commit to a recovery figure that a successful prosecution
would deliver.
